% GENERATED FILE. Edit canonical lessons, then run npm run generate:books.
% canonical-source-hash: fnv1a64:145f46cc92a197b6
% canonical-lessons: RU-C429-administrator, RU-C429-zaselyatsya, RU-C429-vyselyatsya, RU-C429-zalog, RU-C429-prodlevat

\chapter{Checking in and out}
\label{ch:ru-checking-in-and-out}
\hlchaptermodality{429}

\begin{chapteropening}
\textbf{By the end of this chapter:} \emph{I can check in and out: the receptionist, checking in, checking out, the deposit and extending a stay.}

Complete the last lesson of chapter 429: I can check in and out: the receptionist, checking in, checking out, the deposit and extending a stay.
\end{chapteropening}

\section[\texorpdfstring{\ru{администратор} (administrátor)}{администратор (administrátor)}]{\ru{администратор} (administrátor) — a receptionist (at a hotel); an administrator}

\label{lesson:RU-C429-administrator}



\begin{quote}
Before the new one: say the Russian for to spend the night, then the Russian for a booking.
\end{quote}

\subsection*{You'll want to know: \ru{администратор}}
\textbf{\ru{администратор}} (\emph{administrátor}) — \textquotedblleft{}a receptionist (at a hotel); an administrator\textquotedblright{}.

\begin{quote}
\ru{Администратор} \ru{дал} \ru{нам} \ru{ключ} \ru{от} \ru{номера}.
\end{quote}

\begin{quote}
\emph{Administrátor dal nam klyuch ot nómera.}
\end{quote}

\begin{quote}
The receptionist gave us the key to the room.
\end{quote}

Masculine, for a woman too. The desk itself is often called \emph{\ru{ресепшен}}.

\begin{cousinweb}
\emph{\ru{администратор}} is usually traced to Latin \emph{administrātor} 'a manager' (EtymDB lists the Latin word), from \emph{administrāre} 'to attend to, to manage'. EtymDB does not list the Russian word.
\end{cousinweb}

\begin{grammarlens}[title={What you've built}]
One more word for telling, explaining and describing at B1.
\end{grammarlens}

\subsection*{Your turn}
\begin{itemize}
\raggedright
  \item \emph{Say it:} \emph{administrátor}
  \item \emph{Say it:} \emph{administrátor}, once more
  \item \emph{Say it:} a sentence of your own with \emph{administrátor}
  \item \emph{Recall:} say the Russian for to spend the night, then the Russian for a booking, then say \emph{administrátor} again
\end{itemize}

\begin{tcolorbox}[breakable,colback=teal!4,colframe=teal!35!black,title={Before you move on}]
What does \ru{администратор} mean? (\textquotedblleft{}A receptionist (at a hotel); an administrator\textquotedblright{}.) Say it once more.
\end{tcolorbox}
\section[\texorpdfstring{\ru{заселяться} (zaselyátsya)}{заселяться (zaselyátsya)}]{\ru{заселяться} (zaselyátsya) — to check into a hotel; to move in}

\label{lesson:RU-C429-zaselyatsya}



\begin{quote}
Before the new one: say the Russian for a booking, then the Russian for a receptionist.
\end{quote}

\subsection*{You'll want to know: \ru{заселяться}}
\textbf{\ru{заселяться}} (\emph{zaselyátsya}) — \textquotedblleft{}to check into a hotel; to move in\textquotedblright{}.

\begin{quote}
\ru{В} \ru{гостиницу} \ru{можно} \ru{заселяться} \ru{после} \ru{двух} \ru{часов}.
\end{quote}

\begin{quote}
\emph{V gostínitsu mózhno zaselyátsya pósle dvukh chasóv.}
\end{quote}

\begin{quote}
You can check in to the hotel after two o'clock.
\end{quote}

Imperfective; perfective \emph{\ru{заселиться}}. Takes \emph{\ru{в}} + accusative.

\begin{cousinweb}
\emph{\ru{заселяться}} is built from \emph{\ru{за}-}, \emph{\ru{селить}} 'to settle' and \emph{-\ru{ся}}. \emph{\ru{Селить}} goes with \emph{\ru{село}} 'a village', which is inherited from Proto-Slavic \emph{selo} (EtymDB).
\end{cousinweb}

\begin{grammarlens}[title={What you've built}]
One more word for telling, explaining and describing at B1.
\end{grammarlens}

\subsection*{Your turn}
\begin{itemize}
\raggedright
  \item \emph{Say it:} \emph{zaselyátsya}
  \item \emph{Say it:} \emph{zaselyátsya}, once more
  \item \emph{Say it:} a sentence of your own with \emph{zaselyátsya}
  \item \emph{Recall:} say the Russian for a booking, then the Russian for a receptionist, then say \emph{zaselyátsya} again
\end{itemize}

\begin{tcolorbox}[breakable,colback=teal!4,colframe=teal!35!black,title={Before you move on}]
What does \ru{заселяться} mean? (\textquotedblleft{}To check into a hotel; to move in\textquotedblright{}.) Say it once more.
\end{tcolorbox}
\section[\texorpdfstring{\ru{выселяться} (vyselyátsya)}{выселяться (vyselyátsya)}]{\ru{выселяться} (vyselyátsya) — to check out; to move out}

\label{lesson:RU-C429-vyselyatsya}



\begin{quote}
Before the new one: say the Russian for a receptionist, then the Russian for to check into a hotel.
\end{quote}

\subsection*{You'll want to know: \ru{выселяться}}
\textbf{\ru{выселяться}} (\emph{vyselyátsya}) — \textquotedblleft{}to check out; to move out\textquotedblright{}.

\begin{quote}
\ru{Из} \ru{номера} \ru{надо} \ru{выселяться} \ru{до} \ru{двенадцати}.
\end{quote}

\begin{quote}
\emph{Iz nómera nádo vyselyátsya do dvenádtsati.}
\end{quote}

\begin{quote}
You have to check out of the room before twelve.
\end{quote}

Imperfective; the perfective is \emph{\ru{выселиться}}. In hotels people also say \emph{\ru{выезжать}}.

\begin{cousinweb}
\emph{\ru{выселяться}} is built from \emph{\ru{вы}-} 'out', \emph{\ru{селить}} 'to settle' and \emph{-\ru{ся}}. \emph{\ru{Селить}} goes with \emph{\ru{село}} 'a village', which is inherited from Proto-Slavic \emph{selo} (EtymDB).
\end{cousinweb}

\begin{grammarlens}[title={What you've built}]
One more word for telling, explaining and describing at B1.
\end{grammarlens}

\subsection*{Your turn}
\begin{itemize}
\raggedright
  \item \emph{Say it:} \emph{vyselyátsya}
  \item \emph{Say it:} \emph{vyselyátsya}, once more
  \item \emph{Say it:} a sentence of your own with \emph{vyselyátsya}
  \item \emph{Recall:} say the Russian for a receptionist, then the Russian for to check into a hotel, then say \emph{vyselyátsya} again
\end{itemize}

\begin{tcolorbox}[breakable,colback=teal!4,colframe=teal!35!black,title={Before you move on}]
What does \ru{выселяться} mean? (\textquotedblleft{}To check out; to move out\textquotedblright{}.) Say it once more.
\end{tcolorbox}
\section[\texorpdfstring{\ru{залог} (zalóg)}{залог (zalóg)}]{\ru{залог} (zalóg) — a deposit (security)}

\label{lesson:RU-C429-zalog}



\begin{quote}
Before the new one: say the Russian for to check into a hotel, then the Russian for to check out.
\end{quote}

\subsection*{You'll want to know: \ru{залог}}
\textbf{\ru{залог}} (\emph{zalóg}) — \textquotedblleft{}a deposit (security)\textquotedblright{}.

\begin{quote}
\ru{При} \ru{заселении} \ru{мы} \ru{оставили} \ru{залог}.
\end{quote}

\begin{quote}
\emph{Pri zaselénii my ostávili zalóg.}
\end{quote}

\begin{quote}
When we checked in we left a deposit.
\end{quote}

Masculine. \emph{\ru{Под} \ru{залог}} — against a deposit; also on bail.

\begin{cousinweb}
\emph{\ru{залог}} is built from \emph{\ru{за}-} and \emph{-\ru{лог}}, the root of \emph{\ru{положить}} 'to lay': something laid down as a pledge. EtymDB lists the word with the senses 'pledge, deposit, security'.
\end{cousinweb}

\begin{grammarlens}[title={What you've built}]
One more word for telling, explaining and describing at B1.
\end{grammarlens}

\subsection*{Your turn}
\begin{itemize}
\raggedright
  \item \emph{Say it:} \emph{zalóg}
  \item \emph{Say it:} \emph{zalóg}, once more
  \item \emph{Say it:} a sentence of your own with \emph{zalóg}
  \item \emph{Recall:} say the Russian for to check into a hotel, then the Russian for to check out, then say \emph{zalóg} again
\end{itemize}

\begin{tcolorbox}[breakable,colback=teal!4,colframe=teal!35!black,title={Before you move on}]
What does \ru{залог} mean? (\textquotedblleft{}A deposit (security)\textquotedblright{}.) Say it once more.
\end{tcolorbox}
\section[\texorpdfstring{\ru{продлевать} (prodlevát)}{продлевать (prodlevát)}]{\ru{продлевать} (prodlevát) — to extend (a stay, a visa)}

\label{lesson:RU-C429-prodlevat}



\begin{quote}
Before the new one: say the Russian for to check out, then the Russian for a deposit.
\end{quote}

\subsection*{You'll want to know: \ru{продлевать}}
\textbf{\ru{продлевать}} (\emph{prodlevát}) — \textquotedblleft{}to extend (a stay, a visa)\textquotedblright{}.

\begin{quote}
\ru{Каждый} \ru{месяц} \ru{он} \ru{продлевает} \ru{визу}.
\end{quote}

\begin{quote}
\emph{Kázhdyy mésyats on prodleváyet vízu.}
\end{quote}

\begin{quote}
Every month he extends his visa.
\end{quote}

Imperfective; perfective \emph{\ru{продлить}} (\emph{\ru{продлю}, \ru{продлишь}}).

\begin{cousinweb}
\emph{\ru{продлевать}} is built from \emph{\ru{про}-} and \emph{\ru{длить}} 'to make long', which is inherited from Proto-Slavic \emph{dĭliti} (EtymDB).
\end{cousinweb}

\begin{grammarlens}[title={What you've built}]
One more word for telling, explaining and describing at B1.
\end{grammarlens}

\subsection*{Your turn}
\begin{itemize}
\raggedright
  \item \emph{Say it:} \emph{prodlevát}
  \item \emph{Say it:} \emph{prodlevát}, once more
  \item \emph{Say it:} a sentence of your own with \emph{prodlevát}
  \item \emph{Recall:} say the Russian for to check out, then the Russian for a deposit, then say \emph{prodlevát} again
\end{itemize}

\begin{tcolorbox}[breakable,colback=teal!4,colframe=teal!35!black,title={Before you move on}]
What does \ru{продлевать} mean? (\textquotedblleft{}To extend (a stay, a visa)\textquotedblright{}.) Say it once more.
\end{tcolorbox}
